\documentclass[11pt]{article}
\usepackage[margin=1in]{geometry}
\usepackage{enumitem}
\usepackage{hyperref}
\usepackage{graphicx}
\usepackage{array}
\usepackage{multicol}
\usepackage{longtable}
\usepackage{titlesec}
\usepackage{listings}
\usepackage{xcolor}
\usepackage[a4paper, margin=1in]{geometry}
\usepackage{float}      % For [H] placement
\usepackage{caption}
\usepackage{subcaption}
\usepackage{amsmath}
\usepackage{setspace}
\usepackage[margin=1in]{geometry}
\usepackage{enumitem}
\usepackage{hyperref}
\usepackage{graphicx}
\usepackage{array}
\usepackage{multicol}
\usepackage{longtable}
\usepackage{titlesec}

% Custom colors for code
\definecolor{codegreen}{rgb}{0,0.6,0}
\definecolor{codegray}{rgb}{0.5,0.5,0.5}
\definecolor{codepurple}{rgb}{0.58,0,0.82}
\definecolor{backcolour}{rgb}{0.95,0.95,0.92}

\lstset{
    backgroundcolor=\color{backcolour},   
    commentstyle=\color{codegreen},
    keywordstyle=\color{magenta},
    numberstyle=\tiny\color{codegray},
    stringstyle=\color{codepurple},
    basicstyle=\ttfamily\footnotesize,
    breakatwhitespace=false,         
    breaklines=true,                 
    captionpos=b,                    
    keepspaces=true,                 
    numbers=none,                    
    numbersep=5pt,                  
    showspaces=false,                
    showstringspaces=false,
    showtabs=false,                  
    tabsize=2,
    language=Python
}

\begin{document}

%==================================================
\begin{center}
    \large \textbf{Sri Sivasubramaniya Nadar College of Engineering, Chennai} \\
    (An Autonomous Institution Affiliated to Anna University) \\
    \vspace{0.3cm}
\end{center}

\begin{table}[!h]
\renewcommand{\arraystretch}{1.5}
\resizebox{\textwidth}{!}{%
\begin{tabular}{|l|cll|}
\hline
Degree \& Branch & \multicolumn{1}{c|}{B.E. Computer Science \& Engineering} & Semester & VI \\ \hline
Subject Code \& Name & \multicolumn{3}{c|}{UCS2612 -- Machine Learning Algorithms Laboratory} \\ \hline
Academic Year & 2025--2026 (Even) & Batch & 2023--2027 \\ \hline
Due Date & \multicolumn{3}{c|}{29-03-2026} \\ \hline
Name & \multicolumn{3}{c|}{Priya Verma - 3122235001102} \\ \hline

\end{tabular}
}
\end{table}
\begin{center}
    \textbf{\Large Experiment 7}
\end{center}

\begin{center}
    \textbf{Bagging, Boosting, and Stacked Ensemble Models}
\end{center}

%==================================================
\section{Objective}
\begin{itemize}
    \item To understand ensemble learning strategies: Bagging, Boosting, and Stacking.
    \item To implement Bagging and Boosting classifiers.
    \item To build a Stacked Ensemble model using multiple base learners.
    \item To compare ensemble models in terms of accuracy, stability, and generalization.
    \item To analyze the effect of ensemble methods on bias and variance.
\end{itemize}

%==================================================
\section{Dataset Description}
\textbf{Wisconsin Diagnostic Breast Cancer Dataset}
\begin{itemize}
    \item Total samples: 569
    \item Features: 30 numerical attributes
    \item Target classes: Malignant (M) and Benign (B)
\end{itemize}

%==================================================
\section{Implementation Details}
The following code snippets illustrate the core implementation of the classifiers:

\subsection*{Important Libraries to be included}
\begin{lstlisting}[language=Python]
import numpy as np
import pandas as pd
import matplotlib.pyplot as plt
import seaborn as sns
import warnings
warnings.filterwarnings("ignore")
 
from sklearn.datasets import load_breast_cancer
from sklearn.preprocessing import StandardScaler, LabelEncoder
from sklearn.model_selection import (
    train_test_split, KFold, cross_val_score,
    cross_validate, StratifiedKFold
)
from sklearn.tree import DecisionTreeClassifier
from sklearn.ensemble import (
    BaggingClassifier, AdaBoostClassifier,
    GradientBoostingClassifier, StackingClassifier
)
from sklearn.svm import SVC
from sklearn.naive_bayes import GaussianNB
from sklearn.linear_model import LogisticRegression
from sklearn.metrics import (
    accuracy_score, precision_score, recall_score,
    f1_score, confusion_matrix, roc_curve, auc,
    classification_report
)

plt.rcParams.update({
    "font.family": "Arial",
    "font.size": 15,
    "axes.titlesize": 15,
    "axes.titleweight": "bold",
    "axes.labelsize": 15,
    "axes.labelweight": "bold",
    "xtick.labelsize": 13,
    "ytick.labelsize": 13
})
\end{lstlisting}

\subsection*{1. Load and preprocess the dataset.}
\begin{lstlisting}[language=Python]
# Load directly from sklearn (same as UCI Breast Cancer Wisconsin)
data = load_breast_cancer()

X = pd.DataFrame(data.data, columns=data.feature_names)
y = pd.Series(data.target)           # 0 = Malignant, 1 = Benign

# Map to original labels for display
y_labels = pd.Series(data.target_names[data.target])

print(f"Total samples   : {X.shape[0]}")
print(f"Total features  : {X.shape[1]}")
print(f"\nFeature names:\n{list(data.feature_names)}")
print(f"\nClass distribution:")
for cls, name in enumerate(data.target_names):
    count = (y == cls).sum()
    print(f"  {name} ({cls}): {count} samples ({100*count/len(y):.1f}%)")

# Standardize features
scaler = StandardScaler()
X_scaled = scaler.fit_transform(X)
X_scaled = pd.DataFrame(X_scaled, columns=X.columns)

\end{lstlisting}

\subsection*{2. Perform Exploratory Data Analysis (EDA).}
\begin{lstlisting}[language=Python]
print(f"\nBasic Statistics:\n{X.describe().T[['mean','std','min','max']].round(3)}")
print(f"\nMissing values : {X.isnull().sum().sum()}")

fig = plt.figure(figsize=(18, 14))
fig.suptitle("EDA: Wisconsin Breast Cancer Dataset", fontsize=15, fontweight="bold")

# 2a. Class distribution
ax1 = fig.add_subplot(3, 3, 1)
counts = y_labels.value_counts()
bars = ax1.bar(counts.index, counts.values,
               color=["tomato", "steelblue"], edgecolor="black", alpha=0.85)
ax1.set_title("Class Distribution", fontweight="bold")
ax1.set_ylabel("Count")
for bar, val in zip(bars, counts.values):
    ax1.text(bar.get_x() + bar.get_width()/2, bar.get_height(),
             str(val), ha="center", va="bottom", fontweight="bold")
ax1.grid(axis="y", alpha=0.4)

# 2b. Feature correlation heatmap (first 10 features)
ax2 = fig.add_subplot(3, 3, 2)
corr = X.iloc[:, :10].corr()
sns.heatmap(corr, ax=ax2, cmap="coolwarm", center=0,
            annot=False, linewidths=0.3, cbar_kws={"shrink": 0.8})
ax2.set_title("Correlation (First 10 Features)", fontweight="bold")
ax2.tick_params(axis="x", rotation=90, labelsize=7)
ax2.tick_params(axis="y", labelsize=7)

# 2c–2h. Box plots for top 6 features by class separability
top_features = X.columns[:6]
for i, feat in enumerate(top_features):
    ax = fig.add_subplot(3, 3, i + 3)
    for cls, name, color in zip([0, 1], data.target_names, ["tomato", "steelblue"]):
        vals = X.loc[y == cls, feat]
        ax.hist(vals, bins=20, alpha=0.6, color=color, label=name, edgecolor="black", linewidth=0.4)
    ax.set_title(feat[:30], fontsize=8, fontweight="bold")
    ax.set_xlabel("Value", fontsize=7)
    ax.set_ylabel("Count", fontsize=7)
    ax.legend(fontsize=6)
    ax.grid(alpha=0.3)
 
plt.tight_layout()
plt.savefig("eda_analysis.png", dpi=150)
plt.show()
print("Saved: eda_analysis.png")
\end{lstlisting}

\subsection*{3. Split the dataset into training and testing sets (80--20).}
\begin{lstlisting}[language=Python]
# 80-20 stratified split
X_train, X_test, y_train, y_test = train_test_split(
    X_scaled, y, test_size=0.2, random_state=42, stratify=y
)
print(f"Training samples : {X_train.shape[0]}")
print(f"Testing  samples : {X_test.shape[0]}")
print(f"Train class dist : {dict(y_train.value_counts().sort_index())}")
print(f"Test  class dist : {dict(y_test.value_counts().sort_index())}")
\end{lstlisting}

\subsection*{4. Implement a Bagging classifier using a base estimator.}
\begin{lstlisting}[language=Python]
# 5-Fold CV splitter
skf = StratifiedKFold(n_splits=5, shuffle=True, random_state=42)

# Hyperparameter grid
bagging_grid = [
    {"n_estimators": n, "max_samples": ms}
    for n in [10, 50, 100, 200]
    for ms in [0.5, 0.7, 1.0]
]

print("\nTable 1: Bagging Hyperparameter Evaluation")
print(f"{'n_estimators':>14} {'max_samples':>12} {'Avg CV Acc (%)':>16} {'Avg CV F1':>12}")
print("-" * 60)

best_bag_score = 0; best_bag_params = {}; bagging_results = []
for params in bagging_grid:
    clf = BaggingClassifier(
        estimator=DecisionTreeClassifier(random_state=42),
        n_estimators=params["n_estimators"],
        max_samples=params["max_samples"],
        random_state=42, n_jobs=-1
    )
    acc_scores = cross_val_score(clf, X_train, y_train, cv=skf, scoring="accuracy")
    f1_scores  = cross_val_score(clf, X_train, y_train, cv=skf, scoring="f1")
    avg_acc = acc_scores.mean() * 100
    avg_f1  = f1_scores.mean()
    bagging_results.append({**params, "acc": avg_acc, "f1": avg_f1})
    print(f"{params['n_estimators']:>14} {params['max_samples']:>12.1f} "
          f"{avg_acc:>16.2f} {avg_f1:>12.4f}")
    if avg_acc > best_bag_score:
        best_bag_score = avg_acc
        best_bag_params = params
 
print(f"\nBest Bagging Params: {best_bag_params}  |  CV Acc: {best_bag_score:.2f}%")

# Train final Bagging model
bagging_model = BaggingClassifier(
    estimator=DecisionTreeClassifier(random_state=42),
    n_estimators=best_bag_params["n_estimators"],
    max_samples=best_bag_params["max_samples"],
    random_state=42, n_jobs=-1
)
bagging_model.fit(X_train, y_train)
y_pred_bag = bagging_model.predict(X_test)
\end{lstlisting}

\subsection*{5. Implement Boosting classifiers (AdaBoost and/or Gradient Boosting).}
\begin{lstlisting}[language=Python]
# ── 5a. AdaBoost ──────────────────────────────────────────────
print("\n--- AdaBoost ---")
ada_grid = [
    {"n_estimators": n, "learning_rate": lr}
    for n in [50, 100, 200]
    for lr in [0.01, 0.1, 0.5, 1.0]
]
 
print("\nTable 2a: AdaBoost Hyperparameter Evaluation")
print(f"{'n_estimators':>14} {'learning_rate':>14} {'Avg CV Acc (%)':>16} {'Avg CV F1':>12}")
print("-" * 62)

best_ada_score = 0; best_ada_params = {}
for params in ada_grid:
    clf = AdaBoostClassifier(
        n_estimators=params["n_estimators"],
        learning_rate=params["learning_rate"],
        random_state=42
    )
    acc_scores = cross_val_score(clf, X_train, y_train, cv=skf, scoring="accuracy")
    f1_scores  = cross_val_score(clf, X_train, y_train, cv=skf, scoring="f1")
    avg_acc = acc_scores.mean() * 100
    avg_f1  = f1_scores.mean()
    print(f"{params['n_estimators']:>14} {params['learning_rate']:>14.2f} "
          f"{avg_acc:>16.2f} {avg_f1:>12.4f}")
    if avg_acc > best_ada_score:
        best_ada_score = avg_acc
        best_ada_params = params
 
print(f"\nBest AdaBoost Params: {best_ada_params}  |  CV Acc: {best_ada_score:.2f}%")

# ── 5b. Gradient Boosting ─────────────────────────────────────
print("\n--- Gradient Boosting ---")
gb_grid = [
    {"n_estimators": n, "learning_rate": lr, "max_depth": d}
    for n in [50, 100, 200]
    for lr in [0.05, 0.1, 0.2]
    for d in [3, 5]
]
 
print("\nTable 2b: Gradient Boosting Hyperparameter Evaluation")
print(f"{'n_est':>6} {'lr':>6} {'max_d':>7} {'Avg CV Acc (%)':>16} {'Avg CV F1':>12}")
print("-" * 55)
best_gb_score = 0; best_gb_params = {}; boosting_results = []
for params in gb_grid:
    clf = GradientBoostingClassifier(
        n_estimators=params["n_estimators"],
        learning_rate=params["learning_rate"],
        max_depth=params["max_depth"],
        random_state=42
    )
    acc_scores = cross_val_score(clf, X_train, y_train, cv=skf, scoring="accuracy")
    f1_scores  = cross_val_score(clf, X_train, y_train, cv=skf, scoring="f1")
    avg_acc = acc_scores.mean() * 100
    avg_f1  = f1_scores.mean()
    boosting_results.append({**params, "acc": avg_acc, "f1": avg_f1})
    print(f"{params['n_estimators']:>6} {params['learning_rate']:>6.2f} "
          f"{params['max_depth']:>7} {avg_acc:>16.2f} {avg_f1:>12.4f}")
    if avg_acc > best_gb_score:
        best_gb_score = avg_acc
        best_gb_params = params
 
print(f"\nBest Gradient Boosting Params: {best_gb_params}  |  CV Acc: {best_gb_score:.2f}%")

# Train final Boosting models
ada_model = AdaBoostClassifier(
    n_estimators=best_ada_params["n_estimators"],
    learning_rate=best_ada_params["learning_rate"],
    random_state=42
)
ada_model.fit(X_train, y_train)
y_pred_ada = ada_model.predict(X_test)
 
gb_model = GradientBoostingClassifier(
    n_estimators=best_gb_params["n_estimators"],
    learning_rate=best_gb_params["learning_rate"],
    max_depth=best_gb_params["max_depth"],
    random_state=42
)
gb_model.fit(X_train, y_train)
y_pred_gb = gb_model.predict(X_test)
\end{lstlisting}

\subsection*{6. Construct a Stacked Ensemble using multiple base models.}
\begin{lstlisting}[language=Python]
# Base models: SVM, Naive Bayes, Decision Tree
# Meta-learner: Logistic Regression
base_estimators = [
    ("svm",  SVC(kernel="rbf", C=1.0, probability=True, random_state=42)),
    ("nb",   GaussianNB()),
    ("dt",   DecisionTreeClassifier(max_depth=5, random_state=42)),
]
meta_learner = LogisticRegression(C=1.0, max_iter=1000, random_state=42)
 
stacking_model = StackingClassifier(
    estimators=base_estimators,
    final_estimator=meta_learner,
    cv=5,
    stack_method="predict_proba",
    n_jobs=-1
)

\end{lstlisting}

\begin{figure}[H]
    \centering
    \includegraphics[width=1.00\textwidth]{dt.png}
    \caption{Best Parameters for Decision Tree using 5-Fold Cross Validation}
\end{figure}

\subsection*{7. Define hyperparameter search spaces for each ensemble model.}
\begin{lstlisting}[language=Python]
# ── Hyperparameter search space ───────────────────────────────
# We vary: SVM kernel & C, DT max_depth, meta-learner C
stacking_grid = [{"svm_kernel": "rbf",  "svm_C": 0.1,
        "dt_depth":   3,      "meta_C": 0.1,},
    {
        "svm_kernel": "rbf",  "svm_C": 1.0,
        "dt_depth":   3,      "meta_C": 1.0,
    },
    {
        "svm_kernel": "rbf",  "svm_C": 10.0,
        "dt_depth":   5,      "meta_C": 1.0,
    },
    {
        "svm_kernel": "linear", "svm_C": 0.1,
        "dt_depth":   3,        "meta_C": 0.1,
    },
    {
        "svm_kernel": "linear", "svm_C": 1.0,
        "dt_depth":   5,        "meta_C": 1.0,
    },
    {
        "svm_kernel": "linear", "svm_C": 10.0,
        "dt_depth":   5,        "meta_C": 10.0,
    },
    {
        "svm_kernel": "rbf",  "svm_C": 1.0,
        "dt_depth":   None,   "meta_C": 1.0,   # DT fully grown
    },
    {
        "svm_kernel": "rbf",  "svm_C": 5.0,
        "dt_depth":   7,      "meta_C": 0.5,
    },
]
print("\nTable 3: Stacked Ensemble Hyperparameter Evaluation (5-Fold CV)")
print(f"{'SVM kernel':>11} {'SVM C':>7} {'DT depth':>9} {'Meta C':>7} "
      f"{'Base Models':<35} {'Meta Learner':<22} {'CV Acc (%)':>12} {'CV F1':>10}")
print("-" * 115)
best_stack_score = 0
best_stack_params = {}
best_stack_acc_folds = None
best_stack_f1_folds  = None
stacking_results = []
 
for params in stacking_grid:
    # Build base estimators with current hyperparams
    base_estimators_search = [
        ("svm", SVC(
            kernel=params["svm_kernel"],
            C=params["svm_C"],
            probability=True,
            random_state=42
        )),
        ("nb",  GaussianNB()),
        ("dt",  DecisionTreeClassifier(
            max_depth=params["dt_depth"],
            random_state=42
        )),
    ]
    meta = LogisticRegression(
        C=params["meta_C"],
        max_iter=1000,
        random_state=42
    )
    clf = StackingClassifier(
        estimators=base_estimators_search,
        final_estimator=meta,
        cv=5,
        stack_method="predict_proba",
        n_jobs=-1
    )
 
    acc_scores = cross_val_score(clf, X_train, y_train, cv=skf, scoring="accuracy")
    f1_scores  = cross_val_score(clf, X_train, y_train, cv=skf, scoring="f1")
    avg_acc = acc_scores.mean() * 100
    avg_f1  = f1_scores.mean()
 
    depth_str = str(params["dt_depth"]) if params["dt_depth"] else "None"
    base_str  = f"SVM({params['svm_kernel']}) + NB + DT(d={depth_str})"
    meta_str  = f"LR(C={params['meta_C']})"
 
    stacking_results.append({**params, "acc": avg_acc, "f1": avg_f1,
                              "base_str": base_str, "meta_str": meta_str,
                              "acc_folds": acc_scores, "f1_folds": f1_scores})
 
    print(f"{params['svm_kernel']:>11} {params['svm_C']:>7.1f} {depth_str:>9} "
          f"{params['meta_C']:>7.1f} {base_str:<35} {meta_str:<22} "
          f"{avg_acc:>12.2f} {avg_f1:>10.4f}")
 
    if avg_acc > best_stack_score:
        best_stack_score  = avg_acc
        best_stack_params = params
        best_stack_acc_folds = acc_scores
        best_stack_f1_folds  = f1_scores
        best_base_str = base_str
        best_meta_str = meta_str
        
# Print best configuration
print(f"\n{'─'*65}")
print(f"Best Stacking Config:")
print(f"  SVM kernel  : {best_stack_params['svm_kernel']}")
print(f"  SVM C       : {best_stack_params['svm_C']}")
print(f"  DT max_depth: {best_stack_params['dt_depth']}")
print(f"  Meta LR C   : {best_stack_params['meta_C']}")
print(f"  CV Accuracy : {best_stack_score:.2f}%")
print(f"  Per-fold Acc: {[f'{v*100:.2f}%' for v in best_stack_acc_folds]}")
print(f"  Mean ± Std  : {best_stack_acc_folds.mean()*100:.2f}% ± {best_stack_acc_folds.std()*100:.2f}%")
print(f"{'─'*65}")
\end{lstlisting}

\subsection*{8. Select hyperparameters using 5-fold cross-validation.}
\begin{lstlisting}[language=Python]
# ── Train final Stacking model with best hyperparams ─────────
best_base_estimators = [
    ("svm", SVC(
        kernel=best_stack_params["svm_kernel"],
        C=best_stack_params["svm_C"],
        probability=True,
        random_state=42
    )),
    ("nb",  GaussianNB()),
    ("dt",  DecisionTreeClassifier(
        max_depth=best_stack_params["dt_depth"],
        random_state=42
    )),
]
best_meta = LogisticRegression(
    C=best_stack_params["meta_C"],
    max_iter=1000,
    random_state=42
)
stacking_model = StackingClassifier(
    estimators=best_base_estimators,
    final_estimator=best_meta,
    cv=5,
    stack_method="predict_proba",
    n_jobs=-1
)
stacking_model.fit(X_train, y_train)
y_pred_stack = stacking_model.predict(X_test)
 
stack_acc = best_stack_acc_folds
stack_f1  = best_stack_f1_folds
 
print(f"\nFinal Stacking model trained with best hyperparameters.")
print(f"Test set predictions generated.")
\end{lstlisting}

\subsection*{9. Evaluate all models using standard performance metrics.}
\begin{lstlisting}[language=Python]
def get_metrics(y_true, y_pred, name):
    return {
        "Model":     name,
        "Accuracy":  accuracy_score(y_true, y_pred) * 100,
        "Precision": precision_score(y_true, y_pred, zero_division=0),
        "Recall":    recall_score(y_true, y_pred, zero_division=0),
        "F1 Score":  f1_score(y_true, y_pred, zero_division=0),
    }
 
# Use best boosting model (higher of AdaBoost vs GB)
if best_ada_score >= best_gb_score:
    y_pred_boost = y_pred_ada
    boost_name   = f"AdaBoost (n={best_ada_params['n_estimators']}, lr={best_ada_params['learning_rate']})"
else:
    y_pred_boost = y_pred_gb
    boost_name   = f"Gradient Boosting (n={best_gb_params['n_estimators']}, lr={best_gb_params['learning_rate']})"
\end{lstlisting}

\subsection*{10. Compare results and analyze bias--variance behavior.}
\begin{lstlisting}[language=Python]
results_list = [
    get_metrics(y_test, y_pred_bag,   f"Bagging (n={best_bag_params['n_estimators']})"),
    get_metrics(y_test, y_pred_boost,  boost_name),
    get_metrics(y_test, y_pred_stack, "Stacked Ensemble"),
]
results_df = pd.DataFrame(results_list).set_index("Model")
 
print("\nTable 4: Performance Comparison of Ensemble Models")
print(f"{'Model':<45} {'Accuracy (%)':>13} {'Precision':>11} {'Recall':>8} {'F1 Score':>10}")
print("-" * 93)
for model, row in results_df.iterrows():
    print(f"{model:<45} {row['Accuracy']:>13.2f} {row['Precision']:>11.4f} "
          f"{row['Recall']:>8.4f} {row['F1 Score']:>10.4f}")
          
\end{lstlisting}

%==================================================
% -------------------- SECTION --------------------
\section{Visualizations}
This section presents the visual analysis of the class distribution, feature correlation analysis, confusion matrix and ROC Curve.

% -------------------- FIGURES --------------------

\begin{figure}[H]
    \centering
    \includegraphics[width=\textwidth]{confusion_matrices.png}
    \caption{Confusion Matrices — Ensemble Models}
\end{figure}

\begin{figure}[H]
    \centering
    \includegraphics[width=\textwidth]{eda_analysis.png}
    \caption{EDA: Wisconsin Breast Cancer Dataset}
\end{figure}

\begin{figure}[H]
    \centering
    \includegraphics[width=\textwidth]{bagging_hyperparams.png}
    \caption{Bagging Hyperparameter Analysis}
\end{figure}

\begin{figure}[H]
    \centering
    \includegraphics[width=\textwidth]{roc_curves.png}
    \caption{ROC Curves — All Ensemble Models}
\end{figure}

\begin{figure}[H]
    \centering
    \includegraphics[width=\textwidth]{boosting_hyperparams.png}
    \caption{Gradient Boosting Hyperparameter Analysis}
\end{figure}

\begin{figure}[H]
    \centering
    \includegraphics[width=\textwidth]{metrics_comparison.png}
    \caption{Ensemble Model Performance Comparison}
\end{figure}


\begin{figure}[H]
    \centering
    \includegraphics[width=\textwidth]{bias_variance_cv.png}
    \caption{Bias–Variance Analysis: CV Accuracy per Fold}
\end{figure}

\begin{figure}[H]
    \centering
    \includegraphics[width=\textwidth]{feature_importances.png}
    \caption{Top 20 Feature Importances — Gradient Boosting}
\end{figure}

\begin{figure}[H]
    \centering
    \includegraphics[width=\textwidth]{adaboost_error_curve.png}
    \caption{AdaBoost: Estimator Error vs Number of Estimators}
\end{figure}

%==================================================
\section{Performance Tables}
%==================================================
\begin{table}[H]
\centering
\renewcommand{\arraystretch}{1.3}
\caption{Bagging Hyperparameter Evaluation}
\begin{tabular}{|l|c|c|c|}
\hline
n\_estimators & max\_samples & Avg CV Accuracy (\%) & Avg CV F1 Score \\
\hline
10 & 0.5 & 94.51 & 0.9556 \\
10 & 0.7 & 95.16 & 0.9611 \\
10 & 1.0 & 94.95 & 0.9591 \\
50 & 0.5 & 94.51 & 0.9554 \\
50 & 0.7 & 95.38 & 0.9631 \\
50 & 1.0 & 95.16 & 0.9611 \\
100 & 0.5 & 94.95 & 0.9591 \\
100 & 0.7 & 95.82 & 0.9665 \\
100 & 1.0 & 95.16 & 0.9609 \\
200 & 0.5 & 94.73 & 0.9575 \\
200 & 0.7 & 95.82 & 0.9666 \\
200 & 1.0 & 95.16 & 0.9611 \\
\hline
\end{tabular}
\end{table}

%==================================================
\begin{table}[H]
\centering
\renewcommand{\arraystretch}{1.3}
\caption{Boosting Hyperparameter Evaluation}
\begin{tabular}{|c|c|c|c|}
\hline
n\_estimators & learning\_rate & Avg CV Accuracy (\%) & Avg CV F1 Score \\
\hline
50 & 0.01 & 92.53 & 0.9407 \\
50 & 0.10 & 95.38 & 0.9633 \\
50 & 0.50 & 96.70 & 0.9737 \\
50 & 1.00 & 96.70 & 0.9740 \\
100 & 0.01 & 93.41 & 0.9480 \\
100 & 0.10 & 96.04 & 0.9686 \\
100 & 0.50 & 97.36 & 0.9792 \\
100 & 1.00 & 98.02 & 0.9844 \\
200 & 0.01 & 93.63 & 0.9496 \\
200 & 0.10 & 97.14 & 0.9773 \\
200 & 0.50 & 96.92 & 0.9758 \\
200 & 1.00 & 97.80 & 0.9826 \\
 \hline
\end{tabular}
\end{table}

%==================================================
\begin{table}[H]
\centering
\renewcommand{\arraystretch}{1.3}
\caption{Stacked Ensemble Evaluation}
\begin{tabular}{|l|c|c|c|}
\hline
Base Models & Meta Learner & Avg CV Accuracy (\%) & Avg CV F1 Score \\
\hline
SVM(rbf) + NB + DT(d=3) & LR(C=0.1) & 94.51 & 0.9568 \\
SVM(rbf) + NB + DT(d=3) & LR(C=1.0) & 95.82 & 0.9667 \\
SVM(rbf) + NB + DT(d=5) & LR(C=1.0) & 96.70 & 0.9738 \\
SVM(linear) + NB + DT(d=3) & LR(C=0.1) & 96.04 & 0.9689 \\
SVM(linear) + NB + DT(d=5) & LR(C=1.0) & 96.04 & 0.9687 \\ 
SVM(linear) + NB + DT(d=5) & LR(C=10.0) & 95.60 & 0.9649 \\
SVM(rbf) + NB + DT(d=None) & LR(C=1.0) & 96.70 & 0.9735 \\
SVM(rbf) + NB + DT(d=7) & LR(C=0.5) & 96.92 & 0.9753 \\
 \hline
\end{tabular}
\end{table}

%==================================================
\begin{table}[H]
\centering
\renewcommand{\arraystretch}{1.3}
\caption{Performance Comparison of Ensemble Models}
\begin{tabular}{|l|c|c|c|c|}
\hline
Model & Accuracy (\%) & Precision & Recall & F1 Score \\
\hline
Bagging (n=100) & 94.74 & 0.9583 & 0.9583 & 0.9583 \\
Boosting (n=100, lr=1.0) & 95.61 & 0.9467 & 0.9861 & 0.9660 \\
Stacked Ensemble & 98.25 & 0.9730 & 1.0000 & 0.9863 \\
 \hline
\end{tabular}
\end{table}

%==================================================
\section{Observations}

\begin{itemize}

\item \textbf{How does Bagging reduce variance?}
→ Bagging trains 100 Decision Trees on different bootstrap samples
      (max\_samples=0.7) and averages their predictions.
      Individual trees overfit to their sample; averaging cancels out
      random errors, reducing overall variance without increasing bias.

\item \textbf{How does Boosting address model bias?}
 → AdaBoost/Gradient Boosting train estimators sequentially. Each new
      tree focuses on samples the previous ensemble misclassified (higher
      weights in AdaBoost; negative gradient in GB). This iteratively
      corrects systematic errors, reducing bias. Best GB params:
      n\_est=200, lr=0.2, max\_depth=3.
      
\item \textbf{Why does stacking benefit from heterogeneous models?}
→ SVM captures non-linear margin boundaries, Naive Bayes provides
      probabilistic feature independence assumptions, Decision Tree
      captures feature interactions. Each learns a different representation.
      The meta-learner (Logistic Regression) learns the optimal linear
      combination of their outputs, exploiting complementary strengths.


\item \textbf{Which ensemble method performed best and why?}
 → Stacked Ensemble achieved the highest test accuracy of
      98.25\%.

\end{itemize}


\section{Conclusion}
\begin{itemize}

\item Bagging, Boosting, and Stacked Ensemble models were implemented and evaluated.  
\item The experiment demonstrates how different ensemble strategies improve predictive performance and generalization compared to individual models.
\end{itemize}

\end{document}